%قبل از شروع مسئله ترجیحاً چیزی ننویسید. اما اگر خواستید بنویسید، در قالب لیستی (فقرات و شمارش) بنویسید.
\شروع‌مسئله[17]  
% صورت اصلی مسئله را در این‌جا بنویسید.
یک آرایه‌ی \متن‌لاتین{n} تایی داریم که همه‌ی اعداد آن در ابتدا یک هستند. \متن‌لاتین{q} درخواست به شکل زیر داده می شود:

اعداد بازه‌ی 
$[l_i, r_i]$
را در عدد 
$p$
که یک عدد اول است،‌ ضرب کنید. هم چنین می‌دانیم که 
$1 \le l_i \le r_i \le n$
و
$2 \le p \le n$
برقرار است.

حال پس از اعمال \متن‌لاتین{q} درخواست، عدد \متن‌لاتین{i}-ام آرایه به 
$a_i$
تبدیل شده است. در آخر امتیاز آرایه برابر مجموع تعداد مقسوم علیه‌های اعداد آرایه است. حال شما با فرض داشتن \متن‌لاتین{n} و \متن‌لاتین{q} درخواست، باید در 
$O(n+q)$
امتیاز آرایه را بدست آورید. هم چنین می‌توانید فرض کنید که هیچ عددی پس از اعمال درخواست ها،‌ بزرگتر از عدد ثابت 
$C$
نخواهد شد.

\solution{
پاسخ. اگر بدانیم هر خانه از آٰرایه، قرار است در 
${p_i}^{x_i}$
ضرب شود که 
$p_i$
برابر با i امین عدد اول است و توان این عامل نیز برابر
$x_i$
باشد، تعداد مقسوم علیه‌های عددی که در این خانه نوشته‌ می‌شود برابر است با:‌ 
$(x_1+1)(x_2+1)(x_3+1)...$.
درنتیجه برای هر خانه می‌خواهیم بدست بیاوریم چه اعداد اولی توانی بیشتر از 0 دارند. 

حال اگر یک درخواست را در نظر بگیرید، عوامل اول یک توان را در یک بازه به علاوه‌ی یک می‌کند. پس باید روشی بیابیم که بتوانیم به وسیله‌ی آن به گونه‌ای این کار را انجام دهیم. برای همین برای هر خانه از آرایه، دو لیست در نظر میگیریم. این دو لیست را برای خانه‌ی i ام، 
$start_i$
 و 
$end_i$
می‌نامیم. حال به ازای هر i ، برای درخواست i ام، 
$p_i$ 
را در 
$start_{l_i}$
و
$end_{r_i}$
اضافه می‌کنیم. این کار نشان‌دهنده‌ی این است که در خانه‌ی 
$l_i$
یک ضرب در 
$p_i$
آغاز می‌شود و در خانه‌ی 
$r_i$
نیز تمام می‌شود. 

حال یک آرایه به نام power و به طول  C در نظر می‌گیریم که در ابتدا همه یک است(خانه‌ی i ام آن نشان‌دهنده‌ي توان عدد i به علاوه‌ی یک، در تجزیه یک عدد است. اگر i اول نباشد این مقدار همیشه یک می‌ماند). می‌خواهیم این آرایه را به گونه‌ای آپدیت کنیم که برای هر خانه از جدول، بفهمیم توان های عوامل اول چقدر است. برای این کار به شکل زیر عمل می‌کنیم:

ابتدا یک متغیر به نام cnt را که حاصل ضرب اعداد آرایه است برابر یک قرار می‌دهیم.
به ازای 
$i=1,2,...,n$
کار های زیر را به ترتیب انجام می‌دهیم.  

- ابتدا به ازای هر عضو لیست
$start_i$
مثل 
$p$
باید توان این عامل را زیاد کنیم. برای این کار cnt را تقسیم بر توان کنونی (
$power_p$)
می‌کنیم،‌ 
$power_p$
را یکی زیاد می‌کنیم. cnt را در 
$power_p$
ضرب می‌کنیم.

- عددی که cnt دارد، برابر تعداد مقسوم علیه های عدد i ام آرایه است.

- به ازای هر عضو لیست
$end_i$
مثل
$p$
باید توان این عامل را کم کنیم. برای این کار cnt را تقسیم بر توان کنونی می‌کنیم، توان کنونی را کم می‌کنیم، cnt را در توان جدید ضرب می‌کنیم.(همانند گام یک)

تمام عملیات های ریاضی ما از 
$O(1)$
خواهد بود زیرا اعداد از عددی ثابت که در صورت سوال ذکر شده بیشتر نمی‌شوند. هم چنین به ازای هر درخواست داریم از 
$O(1)$
کار انجام می‌دهیم و کلا هم n خانه را بررسی می‌کنیم. پس اردر این الگوریتم از 
$O(n+q)$
است.

}
%اگر مسئله چند قسمت داشت به این صورت عمل کنید:
%\شروع{زیرمسائل}
%\زیرمسئله[0] % مسئلۀ ۱آ
%\زیرمسئله[0] % مسئلۀ ۱ب
%\زیرمسئله[0] % مسئلۀ ۱پ
%\زیرمسئله[0] % مسئلۀ ۱ت
%\پایان{زیرمسائل}




%%%%%%%%%%%%%%%%%%%%%%%%%%%%%%%%%%%%%%%%%%%%%%%%%%%%%
% برخی دستورات مفید برای تمرین‌ها. %
%%%%%%%%%%%%%%%%%%%%%%%%%%%%%%%%%%%%%%%%%%%%%%%%%%%%%

% فرمول‌های نمونه:
%$f(x) = \Theta(x)$
%$T(x, y) \leq \log(x) + 2^y + \binom{2n}{n}$
% {\tt code\_function}


% لیست موارد بدون شماره‌گذاری:
%\شروع{فقرات}
%    \فقره مورد اول 
%        \شروع{فقرات}
%            \فقره مورد اول داخلی 
%                \شروع{فقرات}
%                    \فقره اولین مورد خیلی داخلی 
%                    \فقره دومین مورد خیلی داخلی 
%                \پایان{فقرات}
%            \فقره مورد دوم داخلی 
%        \پایان{فقرات}
%    \فقره مورد دوم 
%\پایان{فقرات}

% لیست موارد شماره‌گذاری شده:
%\شروع{شمارش}
%    \فقره یک مورد 
%    \فقره یک مورد دیگر 
%    \فقره باز هم یک مورد دیگر
%\پایان{شمارش}

% برای نوشتن یک فرمول چند خطی که یک شماره بخوره. برای حالت‌های دیگر به سایت زیر مراجعه کنید:
% https://www.overleaf.com/learn/latex/Aligning_equations_with_amsmath
%\begin{align} 
%    \begin{split}
%        (x+y)^3 &= (x+y)^2(x+y) \\
%                &= (x^2+2xy+y^2)(x+y) \\
%                &= (x^3+2x^2y+xy^2) + (x^2y+2xy^2+y^3) \\
%                &= x^3+3x^2y+3xy^2+y^3
%    \end{split}                                 
%\end{align}

% نوشتن ماتریس و آرایه دوبعدی
%\begin{align}
%    A = 
%    \begin{bmatrix}
%        A_{11} & A_{21} \\
%        A_{21} & A_{22}
%    \end{bmatrix}
%\end{align}

% اضافه کردن تصویر با شرح. برای اضافه کردن تصویر بدون شرح، فقط دستور درج‌تصویر را استفاده کنید.
%\شروع{شکل}[h!]
%\تنظیم‌ازوسط
%  \درج‌تصویر[پهنا=0.5\پهنای‌متن]{img.jpg}
%\شرح{شرح تصویر}
%\برچسب{تصویر:۱}
%\پایان{شکل}
%\پایان{زیرمسائل}

%کد زیر را ببینید.
%%\lstinputlisting{source_filename.py}
%\begin{latin}
%\begin{lstlisting}
%# Import the modules
%import sys
%import random
%
%ans = True
%
%while ans:
%question = raw_input("Ask the magic 8 ball a question: (press enter to quit) ")
%
%answers = random.randint(1,8)
%
%if question == "":
%sys.exit()
%
%elif answers == 1:
%print "It is certain"
%
%elif answers == 2:
%print "Outlook good"
%
%elif answers == 3:
%print "You may rely on it"
%
%elif answers == 4:
%print "Ask again later"
%
%elif answers == 5:
%print "Concentrate and ask again"
%
%elif answers == 6:
%print "Reply hazy, try again"
%
%elif answers == 7:
%print "My reply is no"
%
%elif answers == 8:
%print "My sources say no"
%\end{lstlisting}
%\end{latin}
